International Competition in Mathematics for
             Universtiy Students
                      in
              Plovdiv, Bulgaria
                    1994
                                                                                1

PROBLEMS AND SOLUTIONS

   First day — July 29, 1994

   Problem 1. (13 points)
   a) Let A be a n × n, n ≥ 2, symmetric, invertible matrix with real
positive elements. Show that zn ≤ n2 − 2n, where zn is the number of zero
elements in A−1 .
   b) How many zero elements are there in the inverse of the n × n matrix

                           1 1 1 1 ... 1
                                                      
                          1 2 2 2 ... 2 
                                               
                          1 2 1 1 ... 1 
                       A=                      ?
                                               
                          1 2 1 2 ... 2 
                                               
                          .................... 
                           1 2 1 2 ... ...

   Solution. Denote by aij and bij the elements of A and A−1 , respectively.
                             n
                             P
Then for k 6= m we have            aki bim = 0 and from the positivity of aij we
                             i=0
conclude that at least one of {bim : i = 1, 2, . . . , n} is positive and at least
one is negative. Hence we have at least two non-zero elements in every
column of A−1 . This proves part a). For part b) all b ij are zero except
b1,1 = 2, bn,n = (−1)n , bi,i+1 = bi+1,i = (−1)i for i = 1, 2, . . . , n − 1.

   Problem 2. (13 points)
   Let f ∈ C 1 (a, b), lim f (x) = +∞, lim f (x) = −∞ and
                      x→a+                  x→b−
f 0 (x) + f 2 (x) ≥ −1 for x ∈ (a, b). Prove that b − a ≥ π and give an example
where b − a = π.
   Solution. From the inequality we get

                    d                       f 0 (x)
                      (arctg f (x) + x) =             +1≥0
                   dx                     1 + f 2 (x)

for x ∈ (a, b). Thus arctg f (x)+x is non-decreasing in the interval and using
                   π           π
the limits we get + a ≤ − + b. Hence b − a ≥ π. One has equality for
                   2           2
f (x) = cotg x, a = 0, b = π.

   Problem 3. (13 points)
2

   Given a set S of 2n − 1, n ∈ N, different irrational numbers. Prove
that there are n different elements x 1 , x2 , . . . , xn ∈ S such that for all non-
negative rational numbers a1 , a2 , . . . , an with a1 + a2 + · · · + an > 0 we have
that a1 x1 + a2 x2 + · · · + an xn is an irrational number.
    Solution. Let I be the set of irrational numbers, Q – the set of rational
numbers, Q+ = Q ∩ [0, ∞). We work by induction. For n = 1 the statement
is trivial. Let it be true for n − 1. We start to prove it for n. From the
induction argument there are n − 1 different elements x 1 , x2 , . . . , xn−1 ∈ S
such that

         a1 x1 + a2 x2 + · · · + an−1 xn−1 ∈ I
(1)
         for all a1 , a2 , . . . , an ∈ Q+ with a1 + a2 + · · · + an−1 > 0.

Denote the other elements of S by xn , xn+1 , . . . , x2n−1 . Assume the state-
ment is not true for n. Then for k = 0, 1, . . . , n − 1 there are r k ∈ Q such
that
       n−1                                                            n−1
              bik xi + ck xn+k = rk for some bik , ck ∈ Q+ ,
       X                                                              X
(2)                                                                         bik + ck > 0.
        i=1                                                           i=1

Also
          n−1                                            n−1
                dk xn+k = R for some dk ∈ Q+ ,
          X                                              X
(3)                                                            dk > 0, R ∈ Q.
          k=0                                            k=0

If in (2) ck = 0 then (2) contradicts (1). Thus ck 6= 0 and without loss of
                                                       n−1
generality one may take ck = 1. In (2) also                  bik > 0 in view of xn+k ∈ I.
                                                        P
                                                       i=1
Replacing (2) in (3) we get

         n−1          n−1                              n−1    n−1
                                          !                             !
         X            X                                X      X
               dk −         bik xi + rk       = R or                dk bik xi ∈ Q,
         k=0          i=1                              i=1    k=0

which contradicts (1) because of the conditions on b 0 s and d0 s.


   Problem 4. (18 points)
   Let α ∈ R \ {0} and suppose that F and G are linear maps (operators)
from Rn into Rn satisfying F ◦ G − G ◦ F = αF .
   a) Show that for all k ∈ N one has F k ◦ G − G ◦ F k = αkF k .
   b) Show that there exists k ≥ 1 such that F k = 0.
                                                                                                  3

   Solution. For a) using the assumptions we have
                                         k                                                   
       Fk ◦ G − G ◦ Fk =                        F k−i+1 ◦ G ◦ F i−1 − F k−i ◦ G ◦ F i =
                                         X

                                         i=1
                                          k
                                               F k−i ◦ (F ◦ G − G ◦ F ) ◦ F i−1 =
                                         X
                                    =
                                         i=1
                                          k
                                               F k−i ◦ αF ◦ F i−1 = αkF k .
                                         X
                                    =
                                         i=1

    b) Consider the linear operator L(F ) = F ◦G−G◦F acting over all n×n
matrices F . It may have at most n2 different eigenvalues. Assuming that
F k 6= 0 for every k we get that L has infinitely many different eigenvalues
αk in view of a) – a contradiction.

   Problem 5. (18 points)
   a) Let f ∈ C[0, b], g ∈ C(R) and let g be periodic with period b. Prove
       Z b
that           f (x)g(nx)dx has a limit as n → ∞ and
           0
                             Z b                            Z b               Z b
                                                        1
                     lim           f (x)g(nx)dx =                 f (x)dx ·         g(x)dx.
                     n→∞ 0                              b    0                 0

   b) Find                                     Z π
                                                         sin x
                                        lim                        dx.
                                        n→∞ 0        1 + 3cos 2 nx
                                        Z b
   Solution. Set kgk1 =                       |g(x)|dx and
                                         0

                 ω(f, t) = sup {|f (x) − f (y)| : x, y ∈ [0, b], |x − y| ≤ t} .

In view of the uniform continuity of f we have ω(f, t) → 0 as t → 0. Using
the periodicity of g we get
       Z b                          n Z bk/n
                                    X
               f (x)g(nx)dx =                         f (x)g(nx)dx
        0                         k=1 b(k−1)/n
            n
            X                Z bk/n                      n Z bk/n
                                                         X
       =          f (bk/n)               g(nx)dx +                        {f (x) − f (bk/n)}g(nx)dx
            k=1               b(k−1)/n                   k=1 b(k−1)/n
               n             b
            1X
                               Z
       =          f (bk/n)     g(x)dx + O(ω(f, b/n)kgk1 )
            n k=1          0
4

         n   bk/n              b
      1X
                Z                      Z
    =                f (x)dx     g(x)dx
      b k=1 b(k−1)/n         0
           n                           Z bk/n                  !Z
        1X          b                                                   b
      +               f (bk/n) −                     f (x)dx                g(x)dx + O(ω(f, b/n)kgk1 )
        b k=1       n                   b(k−1)/n                    0
          Z b             Z b
      1
    =           f (x)dx         g(x)dx + O(ω(f, b/n)kgk1 ).
      b    0               0

This proves a). For b) we set b = π, f (x) = sin x, g(x) = (1 + 3cos 2 x)−1 .
From a) and
                    Z π                         Z π
                                                                                    π
                          sin xdx = 2,                (1 + 3cos 2 x)−1 dx =
                     0                           0                                  2
we get                                 Z π
                                                 sin x
                                 lim                       dx = 1.
                                n→∞ 0        1 + 3cos 2 nx


    Problem 6. (25 points)
    Let f ∈ C 2 [0, N ] and |f 0 (x)| < 1, f 00 (x) > 0 for every x ∈ [0, N ]. Let
0 ≤ m0 < m1 < · · · < mk ≤ N be integers such that ni = f (mi ) are also
integers for i = 0, 1, . . . , k. Denote b i = ni − ni−1 and ai = mi − mi−1 for
i = 1, 2, . . . , k.
    a) Prove that
                                  b1   b2           bk
                          −1 <       <      < ··· <    < 1.
                                  a1   a2           ak
    b) Prove that for every choice of A > 1 there are no more than N/A
indices j such that aj > A.
    c) Prove that k ≤ 3N 2/3 (i.e. there are no more than 3N 2/3 integer
points on the curve y = f (x), x ∈ [0, N ]).
    Solution. a) For i = 1, 2, . . . , k we have

                     bi = f (mi ) − f (mi−1 ) = (mi − mi−1 )f 0 (xi )

                                  bi                            bi
for some xi ∈ (mi−1 , mi ). Hence    = f 0 (xi ) and so −1 <        < 1. From the
                                  ai                            ai
                                                      bi
convexity of f we have that f 0 is increasing and        = f 0 (xi ) < f 0 (xi+1 ) =
                                                      ai
bi+1
     because of xi < mi < xi+1 .
ai+1
                                                                                5

    b) Set SA = {j ∈ {0, 1, . . . , k} : aj > A}. Then
                                      k
                                      X            X
                   N ≥ m k − m0 =           ai ≥          aj > A|SA |
                                      i=1          j∈SA

and hence |SA | < N/A.
    c) All different fractions in (−1, 1) with denominators less or equal A are
no more 2A2 . Using b) we get k < N/A + 2A2 . Put A = N 1/3 in the above
estimate and get k < 3N 2/3 .


    Second day — July 30, 1994

    Problem 1. (14 points)
    Let f ∈ C 1 [a, b], f (a) = 0 and suppose that λ ∈ R, λ > 0, is such that

                                 |f 0 (x)| ≤ λ|f (x)|

for all x ∈ [a, b]. Is it true that f (x) = 0 for all x ∈ [a, b]?
     Solution. Assume that there is y ∈ (a, b] such that f (y) 6= 0. Without
loss of generality we have f (y) > 0. In view of the continuity of f there exists
c ∈ [a, y) such that f (c) = 0 and f (x) > 0 for x ∈ (c, y]. For x ∈ (c, y] we
have |f 0 (x)| ≤ λf (x). This implies that the function g(x) = ln f (x) − λx is
                                              f 0 (x)
not increasing in (c, y] because of g 0 (x) =         −λ ≤ 0. Thus ln f (x)−λx ≥
                                              f (x)
ln f (y) − λy and f (x) ≥ eλx−λy f (y) for x ∈ (c, y]. Thus

                     0 = f (c) = f (c + 0) ≥ eλc−λy f (y) > 0

— a contradiction. Hence one has f (x) = 0 for all x ∈ [a, b].

   Problem 2. (14 points)
   Let f : R2 → R be given by f (x, y) = (x2 − y 2 )e−x −y .
                                                       2  2


   a) Prove that f attains its minimum and its maximum.
                                             ∂f           ∂f
   b) Determine all points (x, y) such that     (x, y) =     (x, y) = 0 and
                                             ∂x           ∂y
determine for which of them f has global or local minimum or maximum.
   Solution. We have f (1, 0) = e−1 , f (0, 1) = −e−1 and te−t ≤ 2e−2 for
                                            2  2
t ≥ 2. Therefore |f (x, y)| ≤ (x2 + y 2 )e−x −y ≤ 2e−2 < e−1 for (x, y) ∈
                                                                        /
M = {(u, v) : u2 + v 2 ≤ 2} and f cannot attain its minimum and its
6

maximum outside M . Part a) follows from the compactness of M and the
                                                            ∂f
continuity of f . Let (x, y) be a point from part b). From     (x, y) =
                    2  2
                                                            ∂x
2x(1 − x2 + y 2 )e−x −y we get

(1)                             x(1 − x2 + y 2 ) = 0.

Similarly

(2)                             y(1 + x2 − y 2 ) = 0.

All solutions (x, y) of the system (1), (2) are (0, 0), (0, 1), (0, −1), (1, 0)
and (−1, 0). One has f (1, 0) = f (−1, 0) = e −1 and f has global maximum
at the points (1, 0) and (−1, 0). One has f (0, 1) = f (0, −1) = −e −1 and
f has global minimum at the points (0, 1) and (0, −1). The point (0, 0)
                                                        2
is not an extrema point because of f (x, 0) = x 2 e−x > 0 if x 6= 0 and
                   2
f (y, 0) = −y 2 e−y < 0 if y 6= 0.

   Problem 3. (14 points)
   Let f be a real-valued function with n + 1 derivatives at each point of
R. Show that for each pair of real numbers a, b, a < b, such that
                                                            !
                      f (b) + f 0 (b) + · · · + f (n) (b)
                   ln                                           =b−a
                      f (a) + f 0 (a) + · · · + f (n) (a)

there is a number c in the open interval (a, b) for which

                                 f (n+1) (c) = f (c).

Note that ln denotes the natural logarithm.
                                                                    
    Solution. Set g(x) = f (x) + f 0 (x) + · · · + f (n) (x) e−x . From the
assumption one get g(a) = g(b). Then there exists       c ∈ (a, b) such
                                                                       that
  0                                        0
g (c) = 0. Replacing in the last equality g (x) = f (n+1)  (x) − f (x) e−x we
finish the proof.

      Problem 4. (18 points)
      Let A be a n × n diagonal matrix with characteristic polynomial

                       (x − c1 )d1 (x − c2 )d2 . . . (x − ck )dk ,

where c1 , c2 , . . . , ck are distinct (which means that c1 appears d1 times on the
diagonal, c2 appears d2 times on the diagonal, etc. and d1 +d2 +· · ·+dk = n).
                                                                                                            7

Let V be the space of all n × n matrices B such that AB = BA. Prove that
the dimension of V is
                             d21 + d22 + · · · + d2k .

     Solution. Set A = (aij )ni,j=1 , B = (bij )ni,j=1 , AB = (xij )ni,j=1 and
BA = (yij )ni,j=1 . Then xij = aii bij and yij = ajj bij . Thus AB = BA is
equivalent to (aii − ajj )bij = 0 for i, j = 1, 2, . . . , n. Therefore bij = 0 if
aii 6= ajj and bij may be arbitrary if aii = ajj . The number of indices (i, j)
for which aii = ajj = cm for some m = 1, 2, . . . , k is d2m . This gives the
desired result.

    Problem 5. (18 points)
    Let x1 , x2 , . . . , xk be vectors of m-dimensional Euclidian space, such that
x1 +x2 +· · ·+xk = 0. Show that there exists a permutation π of the integers
{1, 2, . . . , k} such that

                n                  k
                                                     !1/2
                                                 2
                X                  X
                        xπ(i) ≤          kxi k                 for each n = 1, 2, . . . , k.
                i=1                i=1

Note that k · k denotes the Euclidian norm.
    Solution. We define π inductively. Set π(1) = 1. Assume π is defined
for i = 1, 2, . . . , n and also
                                   n                 2       n
                                                                   kxπ(i) k2 .
                                   X                         X
(1)                                      xπ(i)           ≤
                                   i=1                       i=1

Note (1) is true for n = 1. We choose π(n + 1) in a way that (1) is fulfilled
                                             n
                                             P
with n + 1 instead of n. Set y =                         xπ(i) and A = {1, 2, . . . , k} \ {π(i) : i =
                                             i=1                                                       !
1, 2, . . . , n}. Assume that (y, xr ) > 0 for all r ∈ A. Then                         y,
                                                                                            P
                                                                                                  xr       >0
                                                                                            r∈A
and in view of y +
                             P
                                  xr = 0 one gets −(y, y) > 0, which is impossible.
                            r∈A
Therefore there is r ∈ A such that

(2)                                         (y, xr ) ≤ 0.

Put π(n + 1) = r. Then using (2) and (1) we have

      n+1           2
                        = ky + xr k2 = kyk2 + 2(y, xr ) + kxr k2 ≤ kyk2 + kxr k2 ≤
      X
            xπ(i)
      i=1
8

                        n                                  n+1
                              kxπ(i) k2 + kxr k2 =               kxπ(i) k2 ,
                        X                                  X
                    ≤
                        i=1                                i=1
which verifies (1) for n + 1. Thus we define π for every n = 1, 2, . . . , k.
Finally from (1) we get
                      n             2       n                    k
                                                  kxπ(i) k2 ≤          kxi k2 .
                      X                     X                    X
                            xπ(i)       ≤
                      i=1                   i=1                  i=1



    Problem 6. (22 points)
              ln2 N NX−2
                                  1
    Find lim                              . Note that ln denotes the natural
         N →∞ N          ln k · ln(N − k)
                     k=2
logarithm.
    Solution. Obviously

              ln2 N NX
                     −2
                               1         ln2 N N − 3    3
(1)    AN =                            ≥      · 2    =1− .
                N k=2 ln k · ln(N − k)     N   ln N     N

                                                   1
Take M , 2 ≤ M < N/2. Then using that                      is decreasing in
                                          ln k · ln(N − k)
[2, N/2] and the symmetry with respect to N/2 one get
                                                            
                      M   N −M    N
              ln2 N  X     X−1   X −2 
                                                1
         AN =           +       +                       ≤
                N k=2 k=M +1 k=N −M  ln k · ln(N − k)

              ln2 N              M −1           N − 2M − 1
                                                                                 
            ≤             2                 +                                         ≤
                N           ln 2 · ln(N − 2) ln M · ln(N − M )
               2 M ln N           2M                       ln N      1
                                                                               
              ≤    ·         + 1−                               +O                    .
              ln 2     N           N                       ln M    ln N
               N
                   
Choose M =            + 1 to get
             ln2 N
                     2                ln N             1                                      ln ln N
                                                                                                   
(2) AN ≤ 1 −                                      +O                          ≤ 1+O                         .
                  N ln2 N        ln N − 2 ln ln N    ln N                                      ln N
Estimates (1) and (2) give

                        ln2 N NX−2
                                            1
                      lim                           = 1.
                    N →∞ N         ln k · ln(N − k)
                               k=2
